% ======================================================================
%  UPPGIFT 4: JULKLAPPSFÖRDELNING
% ======================================================================
\section{4 Julklappsfördelning}

\uppgiftsbild{4}{Julklappsfördelning}
{modellera en tilldelning som bipartit matchning, lösa den med flöde — och ange tiden i
väldefinierade storheter.}
{\begin{itemize}\setlength{\itemsep}{1pt}
  \item Modellen: vad är hörn, kanter och kapaciteter — och varför just 1?
  \item Girigt räcker inte: en förbättrande stig kan ta tillbaka det som redan delats ut.
  \item Vad $m$ betyder avgör vad $O(\dots)$ blir.
\end{itemize}}

\begin{frame}{Uppgift 4: Julklappsfördelning}
\begin{infobox}[fontupper=\footnotesize]
En pappa ska ge sina $n$ barn var sin julklapp. Varje barn har skrivit en önskelista. Pappan vill
ge varje barn en julklapp från barnets önskelista, men inte samma julklapp till flera barn.
Konstruera och analysera en effektiv algoritm. Du kan anta att det finns högst $m$ saker på
varje önskelista.
\end{infobox}
\begin{columns}[T]
\begin{column}{0.4\textwidth}
\small
\textbf{Årets instans.} Tomten ska ge övningsassistenterna var sin klapp.\par\smallskip
\footnotesize
\renewcommand{\arraystretch}{1.15}
\begin{tabular}{@{}ll@{}}
\toprule
Emma & RTX 5090, PS5 Pro\\
Elliot & RTX 5090, Xbox Series X\\
Marcus & RTX 5090\\
Lovisa & RTX 5090, PS5 Pro, Switch 2\\
\bottomrule
\end{tabular}\par\smallskip
Alla vill ha Nvidia-kortet. Det finns ett.
\end{column}
\begin{column}{0.57\textwidth}
\centering
\input{gen/match-lists}
\end{column}
\end{columns}
\end{frame}

\begin{frame}{Första försöket: girigt}
\begin{columns}[T]
\begin{column}{0.52\textwidth}
\centering
\begin{tikzpicture}[scale=0.85,pers/.style={rounded corners=3pt,draw=kthnavy,fill=white,line width=0.7pt,
    minimum width=15mm,minimum height=5.5mm,inner sep=1pt,font=\footnotesize},
  sak/.style={rounded corners=3pt,draw=kthblue,fill=kthlight,line width=0.7pt,
    minimum width=21mm,minimum height=5.5mm,inner sep=1pt,font=\footnotesize},
  me/.style={draw=black!30,line width=0.7pt},
  mm/.style={draw=kthblue,line width=2.2pt},
  mp/.style={->,draw=kthgrass,line width=2.4pt},
  mb/.style={->,draw=kthorange,line width=2pt,dash pattern=on 4pt off 2pt}]
  \node[pers] (Emma) at (2.3,1.65) {Emma};
  \node[pers] (Elliot) at (2.3,0.55) {Elliot};
  \node[pers,fill=coral!30] (Marcus) at (2.3,-0.55) {Marcus};
  \node[pers] (Lovisa) at (2.3,-1.65) {Lovisa};
  \node[sak] (RTX) at (6.5,1.65) {RTX 5090};
  \node[sak] (PS5) at (6.5,0.55) {PS5 Pro};
  \node[sak] (Xbox) at (6.5,-0.55) {Xbox Series X};
  \node[sak] (Switch) at (6.5,-1.65) {Switch 2};
  \draw[me] (Emma.east) -- (PS5.west);
  \draw[me] (Elliot.east) -- (RTX.west);
  \draw[me] (Marcus.east) -- (RTX.west);
  \draw[me] (Lovisa.east) -- (RTX.west);
  \draw[me] (Lovisa.east) -- (Switch.west);
  \draw[mm] (Emma.east) -- (RTX.west);
  \draw[mm] (Elliot.east) -- (Xbox.west);
  \draw[mm] (Lovisa.east) -- (PS5.west);
\end{tikzpicture}

\end{column}
\begin{column}{0.45\textwidth}
\small
I listordning, var och en tar första lediga sak:\par\smallskip
Emma tar RTX-kortet. Elliot tar Xboxen. Marcus vill bara ha RTX-kortet, som är borta.
Lovisa tar PS5:an.\par\smallskip
Tre klappar. Marcus får ingenting.\par\smallskip
{\footnotesize\color{black!55}(Marcus har gjort de här bilderna. Han har noterat utfallet.)}
\end{column}
\end{columns}
\vfill
\begin{columns}[T]
\begin{column}{0.49\textwidth}
\begin{missbox}[fontupper=\footnotesize]
”Ge var och en första lediga sak.” Resultatet beror på ordningen, och ett tidigt val kan
behöva ångras.
\end{missbox}
\end{column}
\begin{column}{0.49\textwidth}
\begin{ahabox}[fontupper=\footnotesize]
\textbf{Maximal} (går inte att utöka) $\neq$ \textbf{största} (flest par). Den giriga är maximal.
Största har fyra par.
\end{ahabox}
\end{column}
\end{columns}
\varstrip{Kontrast}{instansen}{algoritmen: girig eller flöde}{att ett lokalt val kan behöva ångras}
\end{frame}

\begin{frame}{Modellen: matchning som flöde}
\begin{columns}[T]
\begin{column}{0.54\textwidth}
\centering
\begin{tikzpicture}[scale=0.8,pers/.style={rounded corners=3pt,draw=kthnavy,fill=white,line width=0.7pt,
    minimum width=15mm,minimum height=5.5mm,inner sep=1pt,font=\footnotesize},
  sak/.style={rounded corners=3pt,draw=kthblue,fill=kthlight,line width=0.7pt,
    minimum width=21mm,minimum height=5.5mm,inner sep=1pt,font=\footnotesize},
  me/.style={draw=black!30,line width=0.7pt},
  mm/.style={draw=kthblue,line width=2.2pt},
  mp/.style={->,draw=kthgrass,line width=2.4pt},
  mb/.style={->,draw=kthorange,line width=2pt,dash pattern=on 4pt off 2pt}]
  \node[fn] (s) at (0,0) {$s$};
  \node[fn] (t) at (8.8,0) {$t$};
  \node[pers] (Emma) at (2.3,1.65) {Emma};
  \node[pers] (Elliot) at (2.3,0.55) {Elliot};
  \node[pers] (Marcus) at (2.3,-0.55) {Marcus};
  \node[pers] (Lovisa) at (2.3,-1.65) {Lovisa};
  \node[sak] (RTX) at (6.5,1.65) {RTX 5090};
  \node[sak] (PS5) at (6.5,0.55) {PS5 Pro};
  \node[sak] (Xbox) at (6.5,-0.55) {Xbox Series X};
  \node[sak] (Switch) at (6.5,-1.65) {Switch 2};
  \draw[me] (Emma.east) -- (RTX.west);
  \draw[me] (Emma.east) -- (PS5.west);
  \draw[me] (Elliot.east) -- (RTX.west);
  \draw[me] (Elliot.east) -- (Xbox.west);
  \draw[me] (Marcus.east) -- (RTX.west);
  \draw[me] (Lovisa.east) -- (RTX.west);
  \draw[me] (Lovisa.east) -- (PS5.west);
  \draw[me] (Lovisa.east) -- (Switch.west);
  \draw[fe] (s) -- (Emma.west);
  \draw[fe] (s) -- (Elliot.west);
  \draw[fe] (s) -- (Marcus.west);
  \draw[fe] (s) -- (Lovisa.west);
  \draw[fe] (RTX.east) -- (t);
  \draw[fe] (PS5.east) -- (t);
  \draw[fe] (Xbox.east) -- (t);
  \draw[fe] (Switch.east) -- (t);
\end{tikzpicture}
\par
{\scriptsize alla kapaciteter 1}
\end{column}
\begin{column}{0.43\textwidth}
\small
\begin{itemize}\setlength{\itemsep}{3pt}
  \item $s\to$ barn, kapacitet 1: ingen får två klappar.
  \item barn $\to$ sak om saken står på listan.
  \item sak $\to t$, kapacitet 1: ingen sak delas ut två gånger.
  \item Heltalskapaciteter $\Rightarrow$ Ford--Fulkerson ger ett \emph{heltalsflöde}. Varje
        kant bär 0 eller 1, och kanterna i mitten med flöde 1 är matchningen.
  \item Maxflöde = största matchning.
\end{itemize}
\end{column}
\end{columns}
\end{frame}

\begin{frame}{Ford--Fulkerson (med BFS) på julklapparna}
\centering
\only<1>{%
\begin{tikzpicture}[scale=0.95,pers/.style={rounded corners=3pt,draw=kthnavy,fill=white,line width=0.7pt,
    minimum width=15mm,minimum height=5.5mm,inner sep=1pt,font=\footnotesize},
  sak/.style={rounded corners=3pt,draw=kthblue,fill=kthlight,line width=0.7pt,
    minimum width=21mm,minimum height=5.5mm,inner sep=1pt,font=\footnotesize},
  me/.style={draw=black!30,line width=0.7pt},
  mm/.style={draw=kthblue,line width=2.2pt},
  mp/.style={->,draw=kthgrass,line width=2.4pt},
  mb/.style={->,draw=kthorange,line width=2pt,dash pattern=on 4pt off 2pt}]
\useasboundingbox (-0.4,-2.1) rectangle (9.2,2.1);
  \node[fn] (s) at (0,0) {$s$};
  \node[fn] (t) at (8.8,0) {$t$};
  \node[pers] (Emma) at (2.3,1.65) {Emma};
  \node[pers] (Elliot) at (2.3,0.55) {Elliot};
  \node[pers] (Marcus) at (2.3,-0.55) {Marcus};
  \node[pers] (Lovisa) at (2.3,-1.65) {Lovisa};
  \node[sak] (RTX) at (6.5,1.65) {RTX 5090};
  \node[sak] (PS5) at (6.5,0.55) {PS5 Pro};
  \node[sak] (Xbox) at (6.5,-0.55) {Xbox Series X};
  \node[sak] (Switch) at (6.5,-1.65) {Switch 2};
  \draw[me] (Emma.east) -- (RTX.west);
  \draw[me] (Emma.east) -- (PS5.west);
  \draw[me] (Elliot.east) -- (RTX.west);
  \draw[me] (Elliot.east) -- (Xbox.west);
  \draw[me] (Marcus.east) -- (RTX.west);
  \draw[me] (Lovisa.east) -- (RTX.west);
  \draw[me] (Lovisa.east) -- (PS5.west);
  \draw[me] (Lovisa.east) -- (Switch.west);
  \draw[fe] (s) -- (Emma.west);
  \draw[fe] (s) -- (Elliot.west);
  \draw[fe] (s) -- (Marcus.west);
  \draw[fe] (s) -- (Lovisa.west);
  \draw[fe] (RTX.east) -- (t);
  \draw[fe] (PS5.east) -- (t);
  \draw[fe] (Xbox.east) -- (t);
  \draw[fe] (Switch.east) -- (t);
  \draw[mp] (s) -- (Emma.west);
  \draw[mp] (Emma.east) -- (RTX.west);
  \draw[mp] (RTX.east) -- (t);
\end{tikzpicture}}
\only<2>{%
\begin{tikzpicture}[scale=0.95,pers/.style={rounded corners=3pt,draw=kthnavy,fill=white,line width=0.7pt,
    minimum width=15mm,minimum height=5.5mm,inner sep=1pt,font=\footnotesize},
  sak/.style={rounded corners=3pt,draw=kthblue,fill=kthlight,line width=0.7pt,
    minimum width=21mm,minimum height=5.5mm,inner sep=1pt,font=\footnotesize},
  me/.style={draw=black!30,line width=0.7pt},
  mm/.style={draw=kthblue,line width=2.2pt},
  mp/.style={->,draw=kthgrass,line width=2.4pt},
  mb/.style={->,draw=kthorange,line width=2pt,dash pattern=on 4pt off 2pt}]
\useasboundingbox (-0.4,-2.1) rectangle (9.2,2.1);
  \node[fn] (s) at (0,0) {$s$};
  \node[fn] (t) at (8.8,0) {$t$};
  \node[pers] (Emma) at (2.3,1.65) {Emma};
  \node[pers] (Elliot) at (2.3,0.55) {Elliot};
  \node[pers] (Marcus) at (2.3,-0.55) {Marcus};
  \node[pers] (Lovisa) at (2.3,-1.65) {Lovisa};
  \node[sak] (RTX) at (6.5,1.65) {RTX 5090};
  \node[sak] (PS5) at (6.5,0.55) {PS5 Pro};
  \node[sak] (Xbox) at (6.5,-0.55) {Xbox Series X};
  \node[sak] (Switch) at (6.5,-1.65) {Switch 2};
  \draw[me] (Emma.east) -- (PS5.west);
  \draw[me] (Elliot.east) -- (RTX.west);
  \draw[me] (Elliot.east) -- (Xbox.west);
  \draw[me] (Marcus.east) -- (RTX.west);
  \draw[me] (Lovisa.east) -- (RTX.west);
  \draw[me] (Lovisa.east) -- (PS5.west);
  \draw[me] (Lovisa.east) -- (Switch.west);
  \draw[mm,->] (s) -- (Emma.west);
  \draw[fe] (s) -- (Elliot.west);
  \draw[fe] (s) -- (Marcus.west);
  \draw[fe] (s) -- (Lovisa.west);
  \draw[mm,->] (RTX.east) -- (t);
  \draw[fe] (PS5.east) -- (t);
  \draw[fe] (Xbox.east) -- (t);
  \draw[fe] (Switch.east) -- (t);
  \draw[mm] (Emma.east) -- (RTX.west);
  \draw[mp] (s) -- (Elliot.west);
  \draw[mp] (Elliot.east) -- (Xbox.west);
  \draw[mp] (Xbox.east) -- (t);
\end{tikzpicture}}
\only<3>{%
\begin{tikzpicture}[scale=0.95,pers/.style={rounded corners=3pt,draw=kthnavy,fill=white,line width=0.7pt,
    minimum width=15mm,minimum height=5.5mm,inner sep=1pt,font=\footnotesize},
  sak/.style={rounded corners=3pt,draw=kthblue,fill=kthlight,line width=0.7pt,
    minimum width=21mm,minimum height=5.5mm,inner sep=1pt,font=\footnotesize},
  me/.style={draw=black!30,line width=0.7pt},
  mm/.style={draw=kthblue,line width=2.2pt},
  mp/.style={->,draw=kthgrass,line width=2.4pt},
  mb/.style={->,draw=kthorange,line width=2pt,dash pattern=on 4pt off 2pt}]
\useasboundingbox (-0.4,-2.1) rectangle (9.2,2.1);
  \node[fn] (s) at (0,0) {$s$};
  \node[fn] (t) at (8.8,0) {$t$};
  \node[pers] (Emma) at (2.3,1.65) {Emma};
  \node[pers] (Elliot) at (2.3,0.55) {Elliot};
  \node[pers] (Marcus) at (2.3,-0.55) {Marcus};
  \node[pers] (Lovisa) at (2.3,-1.65) {Lovisa};
  \node[sak] (RTX) at (6.5,1.65) {RTX 5090};
  \node[sak] (PS5) at (6.5,0.55) {PS5 Pro};
  \node[sak] (Xbox) at (6.5,-0.55) {Xbox Series X};
  \node[sak] (Switch) at (6.5,-1.65) {Switch 2};
  \draw[me] (Emma.east) -- (PS5.west);
  \draw[me] (Elliot.east) -- (RTX.west);
  \draw[me] (Marcus.east) -- (RTX.west);
  \draw[me] (Lovisa.east) -- (RTX.west);
  \draw[me] (Lovisa.east) -- (PS5.west);
  \draw[me] (Lovisa.east) -- (Switch.west);
  \draw[mm,->] (s) -- (Emma.west);
  \draw[mm,->] (s) -- (Elliot.west);
  \draw[fe] (s) -- (Marcus.west);
  \draw[fe] (s) -- (Lovisa.west);
  \draw[mm,->] (RTX.east) -- (t);
  \draw[fe] (PS5.east) -- (t);
  \draw[mm,->] (Xbox.east) -- (t);
  \draw[fe] (Switch.east) -- (t);
  \draw[mm] (Emma.east) -- (RTX.west);
  \draw[mm] (Elliot.east) -- (Xbox.west);
  \draw[mp] (s) -- (Lovisa.west);
  \draw[mp] (Lovisa.east) -- (PS5.west);
  \draw[mp] (PS5.east) -- (t);
\end{tikzpicture}}
\only<4>{%
\begin{tikzpicture}[scale=0.95,pers/.style={rounded corners=3pt,draw=kthnavy,fill=white,line width=0.7pt,
    minimum width=15mm,minimum height=5.5mm,inner sep=1pt,font=\footnotesize},
  sak/.style={rounded corners=3pt,draw=kthblue,fill=kthlight,line width=0.7pt,
    minimum width=21mm,minimum height=5.5mm,inner sep=1pt,font=\footnotesize},
  me/.style={draw=black!30,line width=0.7pt},
  mm/.style={draw=kthblue,line width=2.2pt},
  mp/.style={->,draw=kthgrass,line width=2.4pt},
  mb/.style={->,draw=kthorange,line width=2pt,dash pattern=on 4pt off 2pt}]
\useasboundingbox (-0.4,-2.1) rectangle (9.2,2.1);
  \node[fn] (s) at (0,0) {$s$};
  \node[fn] (t) at (8.8,0) {$t$};
  \node[pers] (Emma) at (2.3,1.65) {Emma};
  \node[pers] (Elliot) at (2.3,0.55) {Elliot};
  \node[pers] (Marcus) at (2.3,-0.55) {Marcus};
  \node[pers] (Lovisa) at (2.3,-1.65) {Lovisa};
  \node[sak] (RTX) at (6.5,1.65) {RTX 5090};
  \node[sak] (PS5) at (6.5,0.55) {PS5 Pro};
  \node[sak] (Xbox) at (6.5,-0.55) {Xbox Series X};
  \node[sak] (Switch) at (6.5,-1.65) {Switch 2};
  \draw[me] (Emma.east) -- (PS5.west);
  \draw[me] (Elliot.east) -- (RTX.west);
  \draw[me] (Marcus.east) -- (RTX.west);
  \draw[me] (Lovisa.east) -- (RTX.west);
  \draw[me] (Lovisa.east) -- (Switch.west);
  \draw[mm,->] (s) -- (Emma.west);
  \draw[mm,->] (s) -- (Elliot.west);
  \draw[fe] (s) -- (Marcus.west);
  \draw[mm,->] (s) -- (Lovisa.west);
  \draw[mm,->] (RTX.east) -- (t);
  \draw[mm,->] (PS5.east) -- (t);
  \draw[mm,->] (Xbox.east) -- (t);
  \draw[fe] (Switch.east) -- (t);
  \draw[mm] (Emma.east) -- (RTX.west);
  \draw[mm] (Elliot.east) -- (Xbox.west);
  \draw[mm] (Lovisa.east) -- (PS5.west);
  \draw[mp] (s) -- (Marcus.west);
  \draw[mp] (Marcus.east) -- (RTX.west);
  \draw[mb] (RTX.west) -- (Emma.east);
  \draw[mp] (Emma.east) -- (PS5.west);
  \draw[mb] (PS5.west) -- (Lovisa.east);
  \draw[mp] (Lovisa.east) -- (Switch.west);
  \draw[mp] (Switch.east) -- (t);
\end{tikzpicture}}
\only<5>{%
\begin{tikzpicture}[scale=0.95,pers/.style={rounded corners=3pt,draw=kthnavy,fill=white,line width=0.7pt,
    minimum width=15mm,minimum height=5.5mm,inner sep=1pt,font=\footnotesize},
  sak/.style={rounded corners=3pt,draw=kthblue,fill=kthlight,line width=0.7pt,
    minimum width=21mm,minimum height=5.5mm,inner sep=1pt,font=\footnotesize},
  me/.style={draw=black!30,line width=0.7pt},
  mm/.style={draw=kthblue,line width=2.2pt},
  mp/.style={->,draw=kthgrass,line width=2.4pt},
  mb/.style={->,draw=kthorange,line width=2pt,dash pattern=on 4pt off 2pt}]
\useasboundingbox (-0.4,-2.1) rectangle (9.2,2.1);
  \node[fn] (s) at (0,0) {$s$};
  \node[fn] (t) at (8.8,0) {$t$};
  \node[pers] (Emma) at (2.3,1.65) {Emma};
  \node[pers] (Elliot) at (2.3,0.55) {Elliot};
  \node[pers] (Marcus) at (2.3,-0.55) {Marcus};
  \node[pers] (Lovisa) at (2.3,-1.65) {Lovisa};
  \node[sak] (RTX) at (6.5,1.65) {RTX 5090};
  \node[sak] (PS5) at (6.5,0.55) {PS5 Pro};
  \node[sak] (Xbox) at (6.5,-0.55) {Xbox Series X};
  \node[sak] (Switch) at (6.5,-1.65) {Switch 2};
  \draw[me] (Emma.east) -- (RTX.west);
  \draw[me] (Elliot.east) -- (RTX.west);
  \draw[me] (Lovisa.east) -- (RTX.west);
  \draw[me] (Lovisa.east) -- (PS5.west);
  \draw[mm,->] (s) -- (Emma.west);
  \draw[mm,->] (s) -- (Elliot.west);
  \draw[mm,->] (s) -- (Marcus.west);
  \draw[mm,->] (s) -- (Lovisa.west);
  \draw[mm,->] (RTX.east) -- (t);
  \draw[mm,->] (PS5.east) -- (t);
  \draw[mm,->] (Xbox.east) -- (t);
  \draw[mm,->] (Switch.east) -- (t);
  \draw[mm] (Emma.east) -- (PS5.west);
  \draw[mm] (Elliot.east) -- (Xbox.west);
  \draw[mm] (Marcus.east) -- (RTX.west);
  \draw[mm] (Lovisa.east) -- (Switch.west);
\end{tikzpicture}}
\par\smallskip{\scriptsize\textcolor{kthblue}{\textbf{matchat}}\quad\textcolor{kthgrass}{\textbf{förbättrande stig}}\quad\textcolor{kthorange}{\textbf{bakåtkant: ångra}}}
\vfill
\begin{tcolorbox}[enhanced,halign=left,colback=kthlight!50,colframe=kthsky,boxrule=0.4pt,arc=2pt,
  left=4pt,right=4pt,top=1pt,bottom=1pt,fontupper=\small,height=1.05cm,valign=center]
\only<1>{Alla kapaciteter är 1. BFS hittar $s\to$ Emma $\to$ RTX $\to t$. Emma får grafikkortet. Tills vidare.}\only<2>{$s\to$ Elliot $\to$ Xbox $\to t$: RTX-kortet är upptaget, så Elliot tar sitt andra val.}\only<3>{$s\to$ Lovisa $\to$ PS5 $\to t$. Tre av fyra har en klapp. Marcus står kvar tomhänt.}\only<4>{Enda stigen: Marcus tar RTX från Emma, Emma tar PS5 från Lovisa, Lovisa tar Switch 2. Två bakåtkanter (streckade).}\only<5>{Flöde 4 = perfekt matchning. Stigen i matchningstermer: alternerande ny, bort, ny, bort, ny.}
\end{tcolorbox}
\end{frame}


\begin{frame}{När räcker klapparna inte? Halls villkor}
\begin{columns}[T]
\begin{column}{0.5\textwidth}
\centering
{\footnotesize\bfseries\color{plum}Om Emma också bara vill ha RTX-kortet}\par\smallskip
\begin{tikzpicture}[scale=0.78,pers/.style={rounded corners=3pt,draw=kthnavy,fill=white,line width=0.7pt,
    minimum width=15mm,minimum height=5.5mm,inner sep=1pt,font=\footnotesize},
  sak/.style={rounded corners=3pt,draw=kthblue,fill=kthlight,line width=0.7pt,
    minimum width=21mm,minimum height=5.5mm,inner sep=1pt,font=\footnotesize},
  me/.style={draw=black!30,line width=0.7pt},
  mm/.style={draw=kthblue,line width=2.2pt},
  mp/.style={->,draw=kthgrass,line width=2.4pt},
  mb/.style={->,draw=kthorange,line width=2pt,dash pattern=on 4pt off 2pt}]
  \node[pers,draw=plum,line width=1.2pt] (Emma) at (2.3,1.65) {Emma};
  \node[pers] (Elliot) at (2.3,0.55) {Elliot};
  \node[pers,draw=plum,line width=1.2pt] (Marcus) at (2.3,-0.55) {Marcus};
  \node[pers] (Lovisa) at (2.3,-1.65) {Lovisa};
  \node[sak,draw=plum,line width=1.2pt] (RTX) at (6.5,1.65) {RTX 5090};
  \node[sak] (PS5) at (6.5,0.55) {PS5 Pro};
  \node[sak] (Xbox) at (6.5,-0.55) {Xbox Series X};
  \node[sak] (Switch) at (6.5,-1.65) {Switch 2};
  \draw[me] (Emma.east) -- (RTX.west);
  \draw[me] (Elliot.east) -- (RTX.west);
  \draw[me] (Elliot.east) -- (Xbox.west);
  \draw[me] (Marcus.east) -- (RTX.west);
  \draw[me] (Lovisa.east) -- (RTX.west);
  \draw[me] (Lovisa.east) -- (PS5.west);
  \draw[me] (Lovisa.east) -- (Switch.west);
  \begin{scope}[on background layer]
  \fill[plum!15,rounded corners=4pt] ($(Emma.north west)+(-0.12,0.12)$) rectangle ($(Emma.south east)+(0.12,-0.12)$);
  \fill[plum!15,rounded corners=4pt] ($(Marcus.north west)+(-0.12,0.12)$) rectangle ($(Marcus.south east)+(0.12,-0.12)$);
  \fill[plum!15,rounded corners=4pt] ($(RTX.north west)+(-0.12,0.12)$) rectangle ($(RTX.south east)+(0.12,-0.12)$);
  \end{scope}
\end{tikzpicture}

\end{column}
\begin{column}{0.47\textwidth}
\small
Emma och Marcus önskar tillsammans en enda sak. Två personer, ett kort. Maxflödet blir 3.\par\smallskip
Minsnittet bevisar det: $s\to$ Elliot, $s\to$ Lovisa, RTX $\to t$. Kapacitet 3.\par\smallskip
Vill alla bara ha RTX-kortet blir maxflödet 1: minst tre besvikna övningsassistenter.
\end{column}
\end{columns}
\vfill
\begin{ahabox}
\textbf{Halls sats:} alla kan få en klapp $\iff$ varje grupp $S$ av barn önskar tillsammans minst
$|S|$ olika saker. Annars finns en grupp som bevisar det, och minsnittet hittar den.
\end{ahabox}
\varstrip{Separation}{algoritmen}{önskelistorna}{att minsnittet förklarar varför det inte går}
\end{frame}

\begin{frame}{Vad är $m$? Linjär i vad, igen}
\small
Varje varv i Ford--Fulkerson är en grafsökning: $O(|V|+|E|)$, linjärt i flödesgrafens storlek.
Flödet ökar med 1 per varv och är högst $n$, så det blir högst $n$ varv.
\medskip

\footnotesize
\renewcommand{\arraystretch}{1.25}
\begin{tabularx}{\textwidth}{@{}Xll@{}}
\toprule
\textbf{Om $m$ betyder \dots} & \textbf{kanter i flödesgrafen} & \textbf{tid, $n$ varv}\\
\midrule
högst $m$ saker per lista (Viggo) & $O(nm)$ & $O(n^2m)$\\
totala antalet rader i alla listor & $O(n+m)$ & $O(n(n+m))$\\
antalet olika saker som finns & $O(nm)$ (högst $n$ listor à $m$) & $O(n^2m)$\\
\bottomrule
\end{tabularx}
\medskip
\begin{missbox}
”$O(nm)$, för grafen har $O(nm)$ kanter.” Det är ett varv. Det blir upp till $n$ varv.
\end{missbox}
\varstrip{Separation}{algoritmen}{vad $m$ står för}{att en $O$-formel bara betyder något när variablerna är definierade}
\end{frame}

\begin{frame}{Vad betyder kapaciteterna? En i taget}
\footnotesize
\renewcommand{\arraystretch}{1.3}
\begin{tabularx}{\textwidth}{@{}p{0.27\textwidth}Xc@{}}
\toprule
\textbf{Ändring} & \textbf{Betyder} & \textbf{Rätt modell?}\\
\midrule
$s\to$ barn: 2 & varje barn kan få två klappar & \ja{} (för en annan uppgift)\\
sak $\to t$: $k$ & det finns $k$ exemplar av saken & \ja\\
sak $\to t$: $\infty$ & obegränsat med exemplar: alla får RTX-kortet & \nej{} (Nvidia jublar)\\
$s\to$ barn: $\infty$ & ett barn kan få allt på sin lista & \nej\\
barn $\to$ sak: $\infty$ & ingen skillnad: 1:orna vid $s$ och $t$ begränsar redan & \ja\\
\bottomrule
\end{tabularx}
\vfill
\begin{missbox}
”Alla kapaciteter måste vara 1.” Kanterna i mitten får vara vad som helst $\ge1$. Det är
kanterna vid $s$ och $t$ som uttrycker villkoren.
\end{missbox}
\varstrip{Separation}{grafen}{en kapacitet i taget}{vilken kapacitet som uttrycker vilket villkor}
\end{frame}

{\kaninbg
\begin{frame}{\kh Samma sak utan flöde: alternerande stigar}
\begin{columns}[T]
\begin{column}{0.52\textwidth}
\small
Den sista förbättrande stigen, utan $s$ och $t$:\par\smallskip
\footnotesize
\renewcommand{\arraystretch}{1.15}
\begin{tabular}{@{}lll@{}}
Marcus -- RTX & ny & $+1$\\
RTX -- Emma & bort & $-1$\\
Emma -- PS5 & ny & $+1$\\
PS5 -- Lovisa & bort & $-1$\\
Lovisa -- Switch 2 & ny & $+1$\\
\midrule
& & $+1$ par\\
\end{tabular}\par\smallskip
\small
Börjar och slutar i omatchade hörn, växlar ny/bort. Det är en förbättrande stig i
flödesgrafen, sedd utan $s$ och $t$.
\end{column}
\begin{column}{0.45\textwidth}
\begin{kaninbox}[fontupper=\footnotesize]
\textbf{Berge (1957):} en matchning är störst $\iff$ det inte finns någon alternerande stig
mellan två omatchade hörn.
\end{kaninbox}
\begin{kaninbox}[fontupper=\footnotesize]
\textbf{Hopcroft--Karp (1973):} leta upp många kortaste alternerande stigar, disjunkta, i
samma fas. $O(\sqrt{|V|})$ faser: $O(|E|\sqrt{|V|})$ i stället för $O(|V|\,|E|)$.
\end{kaninbox}
\end{column}
\end{columns}
\end{frame}}

\begin{frame}{Vanliga fel i uppgift 4}
\begin{columns}[T]
\begin{column}{0.49\textwidth}
\begin{missbox}[fontupper=\footnotesize]
”Girigt räcker.” Girigt ger en \emph{maximal} matchning. Här tre par, när fyra går.
\end{missbox}
\begin{missbox}[fontupper=\footnotesize]
”Flödet kan ge Emma och Marcus ett halvt grafikkort var.” Heltalskapaciteter ger ett
heltalsflöde. Ford--Fulkerson halverar inga grafikkort.
\end{missbox}
\end{column}
\begin{column}{0.49\textwidth}
\begin{missbox}[fontupper=\footnotesize]
”Tiden är $O(nm)$.” Ett varv är $O(nm)$. Det blir upp till $n$ varv: $O(n^2m)$.
\end{missbox}
\begin{missbox}[fontupper=\footnotesize]
”Kapacitet $\infty$ på sak $\to t$, för säkerhets skull.” Då får alla samma sak.
\end{missbox}
\end{column}
\end{columns}
\varstrip{Kontrast}{uppgiften}{modell och analys, rätt och nästan rätt}{girig/största, heltal, varv, kapaciteter}
\end{frame}
